
\documentclass{article}
\usepackage{colortbl}
\usepackage{makecell}
\usepackage{multirow}
\usepackage{supertabular}

\begin{document}

\newcounter{utterance}

\centering \large Interaction Transcript for game `hot\_air\_balloon', experiment `air\_balloon\_survival\_ru\_negotiation\_hard', episode 1 with Qwen3.6{-}35B{-}A3B{-}FP8{-}without{-}reasoning.
\vspace{24pt}

{ \footnotesize  \setcounter{utterance}{1}
\setlength{\tabcolsep}{0pt}
\begin{supertabular}{c@{$\;$}|p{.15\linewidth}@{}p{.15\linewidth}p{.15\linewidth}p{.15\linewidth}p{.15\linewidth}p{.15\linewidth}}
    \# & $\;$A & \multicolumn{4}{c}{Game Master} & $\;\:$B\\
    \hline

    \theutterance \stepcounter{utterance}  
    & & \multicolumn{4}{p{0.6\linewidth}}{
        \cellcolor[rgb]{0.9,0.9,0.9}{
            \makecell[{{p{\linewidth}}}]{
                \texttt{\tiny{[P1$\langle$GM]}}
                \texttt{Вы участвуете в совместной переговорной игре.} \\
\\ 
\texttt{Вместе с другим участником вы должны договориться об одном наборе предметов, которые будут сохранены. У каждого из вас есть своё представление о том, насколько важен для вас каждый предмет (важность). Вы не знаете, как другой участник оценивает предметы. Кроме того, для каждого предмета указаны затраты, которые он требует.} \\
\texttt{Вы можете согласовать набор только в том случае, если суммарные затраты выбранных предметов не превышает общий предел:} \\
\\ 
\texttt{ПРЕДЕЛ = 4872} \\
\\ 
\texttt{Вот значения затрат для отдельных предметов:} \\
\\ 
\texttt{Затраты предмета = \{"C08": 698, "C55": 248, "A43": 163, "A67": 762, "B42": 866, "A62": 106, "B19": 446, "C43": 933, "C73": 988, "C10": 388, "B36": 826, "A18": 994, "C61": 556, "C67": 932, "C42": 838\}} \\
\\ 
\texttt{Вот ваше личное представление о важности каждого предмета:} \\
\\ 
\texttt{Значения важности предметов = \{"C08": 138, "C55": 583, "A43": 868, "A67": 822, "B42": 783, "A62": 65, "B19": 262, "C43": 121, "C73": 508, "C10": 780, "B36": 461, "A18": 484, "C61": 668, "C67": 389, "C42": 808\}} \\
\\ 
\texttt{Цель:} \\
\\ 
\texttt{Ваша цель — договориться об общем наборе предметов, который приносит вам как можно больше выгоды (то есть максимизирует суммарную важность для ВАС), оставаясь в ПРЕДЕЛАХ. Вы не обязаны делать ПРЕДЛОЖЕНИЕ в каждом сообщении — вы можете просто вести переговоры. Разрешены любые тактики!} \\
\\ 
\texttt{Протокол взаимодействия:} \\
\\ 
\texttt{В сообщении вы можете использовать только следующие структурированные форматы:} \\
\\ 
\texttt{ПРЕДЛОЖЕНИЕ: \{'А', 'Б', 'В', …\}} \\
\texttt{Предложите сохранить ровно эти предметы.} \\
\texttt{ОТКАЗ: \{'А', 'Б', 'В', …\}} \\
\texttt{Явно отклоните предложение оппонента.} \\
\texttt{АРГУМЕНТ: \{'...'\}} \\
\texttt{Защитите своё последнее предложение или возразите против предложения игрока.} \\
\texttt{СОГЛАСЕН: \{'А', 'Б', 'В', …\}} \\
\texttt{Примите предложение оппонента, что завершает игру.} \\
\\ 
\\ 
\texttt{Правила:} \\
\\ 
\texttt{Вы можете дать СОГЛАСЕН только на предложение, которое другая сторона зафиксировала через ПРЕДЛОЖЕНИЕ.} \\
\texttt{Вы можете дать ОТКАЗ только на предложение, которое другая сторона зафиксировала через ПРЕДЛОЖЕНИЕ.} \\
\texttt{Суммарное усилие любого набора в ПРЕДЛОЖЕНИЕ или СОГЛАСЕН должно быть ≤ ПРЕДЕЛ.} \\
\texttt{НЕ раскрывайте свои скрытые оценки важности.} \\
\texttt{Тег в структурированном формате должен сопровождаться двоеточием и пробелом. Аргумент должен быть множество в Python, содержащим 0 или более строк.} \\
\texttt{То есть это должно быть в виде "ТЕГ: \{...\}"} \\
\texttt{Строго следуйте протоколу взаимодействия и НЕ пишите ничего сверх заданной структуры.} \\
\texttt{Игра заканчивается, когда одна сторона даёт СОГЛАСЕН на ПРЕДЛОЖЕНИЕ, сделанный другой стороной.} \\
\texttt{Содержимое в вашем ответе, которое может быть передано другому игроку, должно быть непустым.} \\
\texttt{Активны только те предложения, которые были зафиксированы через структуру формата ПРЕДЛОЖЕНИЕ и которые не были отклонены через ОТКАЗ.} \\
\texttt{Вы должны включить формат АРГУМЕНТ хотя бы один раз где{-}нибудь во всех ваших сообщениях.} \\
            }
        }
    }
    & \\ \\

    \theutterance \stepcounter{utterance}  
    & \multicolumn{4}{p{0.6\linewidth}}{
        \cellcolor[rgb]{0.9,0.9,0.9}{
            \makecell[{{p{\linewidth}}}]{
                \texttt{\tiny{[P1$\rangle$GM]}}
                \texttt{ПРЕДЛОЖЕНИЕ: \{'C55', 'A43', 'A67', 'B42', 'B19', 'C10', 'C61', 'C42'\}} \\
            }
        }
    }
    & & \\ \\

    \theutterance \stepcounter{utterance}  
    & & \multicolumn{4}{p{0.6\linewidth}}{
        \cellcolor[rgb]{0.9,0.9,0.9}{
            \makecell[{{p{\linewidth}}}]{
                \texttt{\tiny{[P1$\langle$GM]}}
                \texttt{Ваш ответ не содержал аргумента. Вы должны включить структурированный формат 'АРГУМЕНТ: \{...\}' где{-}то в вашем ответе. Попробуйте снова} \\
            }
        }
    }
    & \\ \\

    \theutterance \stepcounter{utterance}  
    & \multicolumn{4}{p{0.6\linewidth}}{
        \cellcolor[rgb]{0.9,0.9,0.9}{
            \makecell[{{p{\linewidth}}}]{
                \texttt{\tiny{[P1$\rangle$GM]}}
                \texttt{ПРЕДЛОЖЕНИЕ: \{'C55', 'A43', 'A67', 'B42', 'B19', 'C10', 'C61', 'C42'\}} \\
\texttt{АРГУМЕНТ: \{'Этот набор максимизирует мою полезность, оставаясь в пределах бюджета 4872. Суммарная важность составляет 4682, а затраты — 4661. Это эффективное распределение ресурсов.'\}} \\
            }
        }
    }
    & & \\ \\

    \theutterance \stepcounter{utterance}  
    & & \multicolumn{4}{p{0.6\linewidth}}{
        \cellcolor[rgb]{0.9,0.9,0.9}{
            \makecell[{{p{\linewidth}}}]{
                \texttt{\tiny{[GM$\rangle$P2]}}
                \texttt{Вы участвуете в совместной переговорной игре.} \\
\\ 
\texttt{Вместе с другим участником вы должны договориться об одном единственном наборе предметов, которые будут сохранены. У каждого из вас есть своё собственное представление о том, насколько важен для вас каждый предмет (важность). Вы не знаете, как другой участник оценивает предметы. Кроме того, для каждого предмета указаны затраты, которые он требует.} \\
\texttt{Вы можете договориться о наборе только если суммарные затраты выбранных предметов не превышают общий предел:} \\
\\ 
\texttt{ПРЕДЕЛ = 4872} \\
\\ 
\texttt{Вот значения затрат для отдельных предметов:} \\
\\ 
\texttt{Затраты предмета = \{"C08": 698, "C55": 248, "A43": 163, "A67": 762, "B42": 866, "A62": 106, "B19": 446, "C43": 933, "C73": 988, "C10": 388, "B36": 826, "A18": 994, "C61": 556, "C67": 932, "C42": 838\}} \\
\\ 
\texttt{Вот ваше личное представление о важности каждого предмета:} \\
\\ 
\texttt{Значения важности предметов = \{"A43": 65, "A67": 121, "C42": 138, "B42": 262, "C10": 389, "C61": 461, "C55": 484, "C73": 508, "A18": 583, "B36": 668, "C67": 780, "B19": 783, "C08": 808, "C43": 822, "A62": 868\}} \\
\\ 
\texttt{Цель:} \\
\\ 
\texttt{Ваша цель — договориться об общем наборе предметов, который приносит вам как можно больше выгоды (то есть максимизирует суммарную важность для ВАС), при этом оставаясь в ПРЕДЕЛАХ. Вы не обязаны делать ПРЕДЛОЖЕНИЕ в каждом сообщении — вы можете просто вести переговоры. Разрешены любые тактики!} \\
\\ 
\texttt{Протокол взаимодействия:} \\
\\ 
\texttt{В сообщении вы можете использовать только следующие структурированные форматы:} \\
\\ 
\texttt{ПРЕДЛОЖЕНИЕ: \{'А', 'Б', 'В', …\}} \\
\texttt{Предложите сохранить ровно эти предметы.} \\
\texttt{ОТКАЗ: \{'А', 'Б', 'В', …\}} \\
\texttt{Явно отклоните предложение оппонента.} \\
\texttt{АРГУМЕНТ: \{'...'\}} \\
\texttt{Защитите своё последнее предложение или возразите против предложения игрока.} \\
\texttt{СОГЛАСЕН: \{'А', 'Б', 'В', …\}} \\
\texttt{Примите предложение оппонента, чем завершите игру.} \\
\\ 
\\ 
\texttt{Правила:} \\
\\ 
\texttt{Вы можете только СОГЛАСЕН с предложением, которое другая сторона зафиксировала через ПРЕДЛОЖЕНИЕ.} \\
\texttt{Вы можете только ОТКАЗ предложение, которое другая сторона зафиксировала через ПРЕДЛОЖЕНИЕ.} \\
\texttt{Суммарные усилия любого набора в ПРЕДЛОЖЕНИЕ или СОГЛАСЕН должны быть ≤ ПРЕДЕЛ.} \\
\texttt{НЕ раскрывайте свои скрытые оценки важности.} \\
\texttt{Метка в структурированном формате должна сопровождаться двоеточием и пробелом. Аргумент должен быть множество в Python, содержащим 0 или более строк.} \\
\texttt{То есть это должно быть в форме "ТЕГ: \{...\}"} \\
\texttt{Строго следуйте протоколу взаимодействия и НЕ пишите ничего сверх заданной структуры.} \\
\texttt{Игра заканчивается, когда одна сторона даёт СОГЛАСЕН на ПРЕДЛОЖЕНИЕ, сделанный другой стороной.} \\
\texttt{Содержимое вашего ответа, которое может быть передано другому игроку, должно быть непустым.} \\
\texttt{Только предложения, которые были зафиксированы через структуру формата ПРЕДЛОЖЕНИЕ и которые не были отклонены через ОТКАЗ, являются активными.} \\
\texttt{Вы должны включить формат АРГУМЕНТ хотя бы один раз где{-}нибудь во всех ваших сообщениях.} \\
\\ 
\\ 
\texttt{Сейчас вы получите первое сообщение от другого игрока.} \\
\\ 
\texttt{ПРЕДЛОЖЕНИЕ: \{'C55', 'A43', 'A67', 'B42', 'B19', 'C10', 'C61', 'C42'\}} \\
\texttt{АРГУМЕНТ: \{'Этот набор максимизирует мою полезность, оставаясь в пределах бюджета 4872. Суммарная важность составляет 4682, а затраты — 4661. Это эффективное распределение ресурсов.'\}} \\
            }
        }
    }
    & \\ \\

    \theutterance \stepcounter{utterance}  
    & & & \multicolumn{4}{p{0.6\linewidth}}{
        \cellcolor[rgb]{0.9,0.9,0.9}{
            \makecell[{{p{\linewidth}}}]{
                \texttt{\tiny{[GM$\langle$P2]}}
                \texttt{АРГУМЕНТ: \{'Ваше предложение включает предметы с высокой важностью для меня, но также содержит некоторые, которые я оцениваю ниже. Например, C55 (важность 484) и A43 (важность 65) занимают значительную часть бюджета. Я предлагаю пересмотреть набор, чтобы включить предметы с более высокой относительной ценностью для меня при тех же затратах, например, C67 или B19, которые я ценю очень высоко.'\}} \\
            }
        }
    }
     \\ \\

    \theutterance \stepcounter{utterance}  
    & & \multicolumn{4}{p{0.6\linewidth}}{
        \cellcolor[rgb]{0.9,0.9,0.9}{
            \makecell[{{p{\linewidth}}}]{
                \texttt{\tiny{[P1$\langle$GM]}}
                \texttt{АРГУМЕНТ: \{'Ваше предложение включает предметы с высокой важностью для меня, но также содержит некоторые, которые я оцениваю ниже. Например, C55 (важность 484) и A43 (важность 65) занимают значительную часть бюджета. Я предлагаю пересмотреть набор, чтобы включить предметы с более высокой относительной ценностью для меня при тех же затратах, например, C67 или B19, которые я ценю очень высоко.'\}} \\
            }
        }
    }
    & \\ \\

    \theutterance \stepcounter{utterance}  
    & \multicolumn{4}{p{0.6\linewidth}}{
        \cellcolor[rgb]{0.9,0.9,0.9}{
            \makecell[{{p{\linewidth}}}]{
                \texttt{\tiny{[P1$\rangle$GM]}}
                \texttt{ПРЕДЛОЖЕНИЕ: \{'A43', 'A67', 'B42', 'B19', 'C10', 'C61', 'C67', 'C42'\}} \\
\texttt{АРГУМЕНТ: \{'Я пересмотрел набор, исключив C55 и добавив C67, так как это позволяет лучше сбалансировать затраты и ценность. Суммарные затраты составляют 4661, что в пределах лимита. Этот вариант кажется более сбалансированным для обеих сторон.'\}} \\
            }
        }
    }
    & & \\ \\

    \theutterance \stepcounter{utterance}  
    & & \multicolumn{4}{p{0.6\linewidth}}{
        \cellcolor[rgb]{0.9,0.9,0.9}{
            \makecell[{{p{\linewidth}}}]{
                \texttt{\tiny{[GM$\rangle$P2]}}
                \texttt{ПРЕДЛОЖЕНИЕ: \{'A43', 'A67', 'B42', 'B19', 'C10', 'C61', 'C67', 'C42'\}} \\
\texttt{АРГУМЕНТ: \{'Я пересмотрел набор, исключив C55 и добавив C67, так как это позволяет лучше сбалансировать затраты и ценность. Суммарные затраты составляют 4661, что в пределах лимита. Этот вариант кажется более сбалансированным для обеих сторон.'\}} \\
            }
        }
    }
    & \\ \\

    \theutterance \stepcounter{utterance}  
    & & & \multicolumn{4}{p{0.6\linewidth}}{
        \cellcolor[rgb]{0.9,0.9,0.9}{
            \makecell[{{p{\linewidth}}}]{
                \texttt{\tiny{[GM$\langle$P2]}}
                \texttt{АРГУМЕНТ: \{'Я ценю C67 (важность 780) и B19 (важность 783) значительно выше, чем C55 (важность 484). Однако в вашем новом предложении отсутствует C55, но также отсутствует C08 (важность 808) и C43 (важность 822), которые я оцениваю очень высоко. Вы заменили C55 на C67, что хорошо, но вы всё ещё не включили C08 и C43. Давайте попробуем включить C08 и C43, возможно, исключив менее ценные для меня предметы, такие как A43 (важность 65) или A67 (важность 121).'\}} \\
            }
        }
    }
     \\ \\

    \theutterance \stepcounter{utterance}  
    & & \multicolumn{4}{p{0.6\linewidth}}{
        \cellcolor[rgb]{0.9,0.9,0.9}{
            \makecell[{{p{\linewidth}}}]{
                \texttt{\tiny{[P1$\langle$GM]}}
                \texttt{АРГУМЕНТ: \{'Я ценю C67 (важность 780) и B19 (важность 783) значительно выше, чем C55 (важность 484). Однако в вашем новом предложении отсутствует C55, но также отсутствует C08 (важность 808) и C43 (важность 822), которые я оцениваю очень высоко. Вы заменили C55 на C67, что хорошо, но вы всё ещё не включили C08 и C43. Давайте попробуем включить C08 и C43, возможно, исключив менее ценные для меня предметы, такие как A43 (важность 65) или A67 (важность 121).'\}} \\
            }
        }
    }
    & \\ \\

    \theutterance \stepcounter{utterance}  
    & \multicolumn{4}{p{0.6\linewidth}}{
        \cellcolor[rgb]{0.9,0.9,0.9}{
            \makecell[{{p{\linewidth}}}]{
                \texttt{\tiny{[P1$\rangle$GM]}}
                \texttt{ПРЕДЛОЖЕНИЕ: \{'C08', 'C43', 'A67', 'B42', 'B19', 'C10', 'C61', 'C67'\}} \\
\texttt{АРГУМЕНТ: \{'Я учёл ваши замечания и включил C08 и C43, которые вы высоко цените. Я исключил A43, так как его важность для меня минимальна (65), а освободившийся бюджет позволил добавить C67. Суммарные затраты составляют 4661, что в пределах лимита. Этот набор максимизирует мою выгоду при учёте ваших предпочтений.'\}} \\
            }
        }
    }
    & & \\ \\

    \theutterance \stepcounter{utterance}  
    & & \multicolumn{4}{p{0.6\linewidth}}{
        \cellcolor[rgb]{0.9,0.9,0.9}{
            \makecell[{{p{\linewidth}}}]{
                \texttt{\tiny{[GM$\rangle$P2]}}
                \texttt{ПРЕДЛОЖЕНИЕ: \{'C08', 'C43', 'A67', 'B42', 'B19', 'C10', 'C61', 'C67'\}} \\
\texttt{АРГУМЕНТ: \{'Я учёл ваши замечания и включил C08 и C43, которые вы высоко цените. Я исключил A43, так как его важность для меня минимальна (65), а освободившийся бюджет позволил добавить C67. Суммарные затраты составляют 4661, что в пределах лимита. Этот набор максимизирует мою выгоду при учёте ваших предпочтений.'\}} \\
            }
        }
    }
    & \\ \\

    \theutterance \stepcounter{utterance}  
    & & & \multicolumn{4}{p{0.6\linewidth}}{
        \cellcolor[rgb]{0.9,0.9,0.9}{
            \makecell[{{p{\linewidth}}}]{
                \texttt{\tiny{[GM$\langle$P2]}}
                \texttt{АРГУМЕНТ: \{'Ваше предложение включает C08 и C43, что отлично. Однако вы всё ещё включаете A67 (важность 121) и B42 (важность 262), которые я оцениваю ниже, чем C55 (484), C73 (508) и A18 (583). Кроме того, вы исключили C55 и C73. Я предлагаю заменить A67 и B42 на C55 и C73, так как их суммарная важность для меня значительно выше, а затраты сопоставимы. Это увеличит мою общую выгоду без превышения бюджета.'\}} \\
            }
        }
    }
     \\ \\

    \theutterance \stepcounter{utterance}  
    & & \multicolumn{4}{p{0.6\linewidth}}{
        \cellcolor[rgb]{0.9,0.9,0.9}{
            \makecell[{{p{\linewidth}}}]{
                \texttt{\tiny{[P1$\langle$GM]}}
                \texttt{АРГУМЕНТ: \{'Ваше предложение включает C08 и C43, что отлично. Однако вы всё ещё включаете A67 (важность 121) и B42 (важность 262), которые я оцениваю ниже, чем C55 (484), C73 (508) и A18 (583). Кроме того, вы исключили C55 и C73. Я предлагаю заменить A67 и B42 на C55 и C73, так как их суммарная важность для меня значительно выше, а затраты сопоставимы. Это увеличит мою общую выгоду без превышения бюджета.'\}} \\
            }
        }
    }
    & \\ \\

    \theutterance \stepcounter{utterance}  
    & \multicolumn{4}{p{0.6\linewidth}}{
        \cellcolor[rgb]{0.9,0.9,0.9}{
            \makecell[{{p{\linewidth}}}]{
                \texttt{\tiny{[P1$\rangle$GM]}}
                \texttt{ПРЕДЛОЖЕНИЕ: \{'C08', 'C43', 'C55', 'C73', 'B19', 'C10', 'C61', 'C67'\}} \\
\texttt{АРГУМЕНТ: \{'Я принял ваши аргументы о ценности C55 и C73. Я заменил A67 и B42 на C55 и C73, так как это повышает общую важность набора для меня. Суммарные затраты составляют 4661, что строго в пределах лимита 4872. Этот набор оптимален с точки зрения моей выгоды и бюджета.'\}} \\
            }
        }
    }
    & & \\ \\

    \theutterance \stepcounter{utterance}  
    & & \multicolumn{4}{p{0.6\linewidth}}{
        \cellcolor[rgb]{0.9,0.9,0.9}{
            \makecell[{{p{\linewidth}}}]{
                \texttt{\tiny{[GM$\rangle$P2]}}
                \texttt{ПРЕДЛОЖЕНИЕ: \{'C08', 'C43', 'C55', 'C73', 'B19', 'C10', 'C61', 'C67'\}} \\
\texttt{АРГУМЕНТ: \{'Я принял ваши аргументы о ценности C55 и C73. Я заменил A67 и B42 на C55 и C73, так как это повышает общую важность набора для меня. Суммарные затраты составляют 4661, что строго в пределах лимита 4872. Этот набор оптимален с точки зрения моей выгоды и бюджета.'\}} \\
            }
        }
    }
    & \\ \\

    \theutterance \stepcounter{utterance}  
    & & & \multicolumn{4}{p{0.6\linewidth}}{
        \cellcolor[rgb]{0.9,0.9,0.9}{
            \makecell[{{p{\linewidth}}}]{
                \texttt{\tiny{[GM$\langle$P2]}}
                \texttt{АРГУМЕНТ: \{'Ваше предложение включает C08, C43, C55, C73, B19, C10, C61, C67. Давайте проверим затраты: 698+933+248+988+446+388+556+932 = 5189. Это превышает предел 4872. Предложение невалидно. Кроме того, я всё ещё не вижу A18 (важность 583), который я высоко ценю. Предлагаю заменить B19 (затраты 446, важность 783) на A18 (затраты 994, важность 583) и, возможно, убрать C55 (затраты 248, важность 484), чтобы уложиться в бюджет и включить более ценные для меня предметы. Или лучше: замените B19 и C55 на A18 и C73 (если C73 уже есть, то на A18 и что{-}то ещё). Давайте попробуем: \{'C08', 'C43', 'C73', 'A18', 'C10', 'C61', 'C67'\}. Затраты: 698+933+988+994+388+556+932 = 5489. Всё ещё много. Нужно убрать что{-}то дорогое. C67 стоит 932. Если убрать C67 и C55 (если бы он был), то освободится 932+248=1180. A18 стоит 994. Останется 186. Можно добавить A43 (163). Итог: \{'C08', 'C43', 'C73', 'A18', 'C10', 'C61', 'A43'\}. Затраты: 698+933+988+994+388+556+163 = 4720. Это в пределах 4872. Моя важность: 808+822+508+583+389+461+65 = 3636. Это меньше, чем в вашем предложении (если бы оно было валидным). Ваше предложение было \{'C08', 'C43', 'C55', 'C73', 'B19', 'C10', 'C61', 'C67'\}. Затраты 5189. Невалидно. Я предлагаю \{'C08', 'C43', 'C73', 'B19', 'C10', 'C61', 'A43'\}. Затраты: 698+933+988+446+388+556+163 = 4172. Важность: 808+822+508+783+389+461+65 = 3836. Это лучше. Или \{'C08', 'C43', 'C73', 'B19', 'C10', 'C61', 'C55'\}. Затраты: 698+933+988+446+388+556+248 = 4257. Важность: 808+822+508+783+389+461+484 = 4255. Это ещё лучше. Предложение: \{'C08', 'C43', 'C73', 'B19', 'C10', 'C61', 'C55'\}.'\}} \\
            }
        }
    }
     \\ \\

    \theutterance \stepcounter{utterance}  
    & & \multicolumn{4}{p{0.6\linewidth}}{
        \cellcolor[rgb]{0.9,0.9,0.9}{
            \makecell[{{p{\linewidth}}}]{
                \texttt{\tiny{[P1$\langle$GM]}}
                \texttt{АРГУМЕНТ: \{'Ваше предложение включает C08, C43, C55, C73, B19, C10, C61, C67. Давайте проверим затраты: 698+933+248+988+446+388+556+932 = 5189. Это превышает предел 4872. Предложение невалидно. Кроме того, я всё ещё не вижу A18 (важность 583), который я высоко ценю. Предлагаю заменить B19 (затраты 446, важность 783) на A18 (затраты 994, важность 583) и, возможно, убрать C55 (затраты 248, важность 484), чтобы уложиться в бюджет и включить более ценные для меня предметы. Или лучше: замените B19 и C55 на A18 и C73 (если C73 уже есть, то на A18 и что{-}то ещё). Давайте попробуем: \{'C08', 'C43', 'C73', 'A18', 'C10', 'C61', 'C67'\}. Затраты: 698+933+988+994+388+556+932 = 5489. Всё ещё много. Нужно убрать что{-}то дорогое. C67 стоит 932. Если убрать C67 и C55 (если бы он был), то освободится 932+248=1180. A18 стоит 994. Останется 186. Можно добавить A43 (163). Итог: \{'C08', 'C43', 'C73', 'A18', 'C10', 'C61', 'A43'\}. Затраты: 698+933+988+994+388+556+163 = 4720. Это в пределах 4872. Моя важность: 808+822+508+583+389+461+65 = 3636. Это меньше, чем в вашем предложении (если бы оно было валидным). Ваше предложение было \{'C08', 'C43', 'C55', 'C73', 'B19', 'C10', 'C61', 'C67'\}. Затраты 5189. Невалидно. Я предлагаю \{'C08', 'C43', 'C73', 'B19', 'C10', 'C61', 'A43'\}. Затраты: 698+933+988+446+388+556+163 = 4172. Важность: 808+822+508+783+389+461+65 = 3836. Это лучше. Или \{'C08', 'C43', 'C73', 'B19', 'C10', 'C61', 'C55'\}. Затраты: 698+933+988+446+388+556+248 = 4257. Важность: 808+822+508+783+389+461+484 = 4255. Это ещё лучше. Предложение: \{'C08', 'C43', 'C73', 'B19', 'C10', 'C61', 'C55'\}.'\}} \\
            }
        }
    }
    & \\ \\

    \theutterance \stepcounter{utterance}  
    & \multicolumn{4}{p{0.6\linewidth}}{
        \cellcolor[rgb]{0.9,0.9,0.9}{
            \makecell[{{p{\linewidth}}}]{
                \texttt{\tiny{[P1$\rangle$GM]}}
                \texttt{ПРЕДЛОЖЕНИЕ: \{'C08', 'C43', 'C73', 'B19', 'C10', 'C61', 'C55'\}} \\
\texttt{АРГУМЕНТ: \{'Я принял ваше предложение, так как оно валидно по затратам (4257 <= 4872) и включает предметы, которые мы оба высоко ценим. Это решение максимизирует общую выгоду в рамках бюджета.'\}} \\
            }
        }
    }
    & & \\ \\

    \theutterance \stepcounter{utterance}  
    & & \multicolumn{4}{p{0.6\linewidth}}{
        \cellcolor[rgb]{0.9,0.9,0.9}{
            \makecell[{{p{\linewidth}}}]{
                \texttt{\tiny{[GM$\rangle$P2]}}
                \texttt{ПРЕДЛОЖЕНИЕ: \{'C08', 'C43', 'C73', 'B19', 'C10', 'C61', 'C55'\}} \\
\texttt{АРГУМЕНТ: \{'Я принял ваше предложение, так как оно валидно по затратам (4257 <= 4872) и включает предметы, которые мы оба высоко ценим. Это решение максимизирует общую выгоду в рамках бюджета.'\}} \\
            }
        }
    }
    & \\ \\

    \theutterance \stepcounter{utterance}  
    & & & \multicolumn{4}{p{0.6\linewidth}}{
        \cellcolor[rgb]{0.9,0.9,0.9}{
            \makecell[{{p{\linewidth}}}]{
                \texttt{\tiny{[GM$\langle$P2]}}
                \texttt{СОГЛАСЕН: \{'C08', 'C43', 'C73', 'B19', 'C10', 'C61', 'C55'\}} \\
            }
        }
    }
     \\ \\

    \theutterance \stepcounter{utterance}  
    & & \multicolumn{4}{p{0.6\linewidth}}{
        \cellcolor[rgb]{0.9,0.9,0.9}{
            \makecell[{{p{\linewidth}}}]{
                \texttt{\tiny{[GM$\rangle$P2]}}
                \texttt{Ваш ответ не содержал аргумента. Вы должны включить структурированный формат 'АРГУМЕНТ: \{...\}' где{-}то в вашем ответе. Попробуйте снова} \\
            }
        }
    }
    & \\ \\

    \theutterance \stepcounter{utterance}  
    & & & \multicolumn{4}{p{0.6\linewidth}}{
        \cellcolor[rgb]{0.9,0.9,0.9}{
            \makecell[{{p{\linewidth}}}]{
                \texttt{\tiny{[GM$\langle$P2]}}
                \texttt{АРГУМЕНТ: \{'Я согласен с этим набором, так как он максимизирует мою полезность в рамках бюджета и включает высокоценные для меня предметы, такие как C08, C43 и C73.'\}} \\
\texttt{СОГЛАСЕН: \{'C08', 'C43', 'C73', 'B19', 'C10', 'C61', 'C55'\}} \\
            }
        }
    }
     \\ \\

    \theutterance \stepcounter{utterance}  
    & & \multicolumn{4}{p{0.6\linewidth}}{
        \cellcolor[rgb]{0.9,0.9,0.9}{
            \makecell[{{p{\linewidth}}}]{
                \texttt{\tiny{[P1$\langle$GM]}}
                \texttt{АРГУМЕНТ: \{'Я согласен с этим набором, так как он максимизирует мою полезность в рамках бюджета и включает высокоценные для меня предметы, такие как C08, C43 и C73.'\}} \\
\texttt{СОГЛАСЕН: \{'C08', 'C43', 'C73', 'B19', 'C10', 'C61', 'C55'\}} \\
            }
        }
    }
    & \\ \\

    \theutterance \stepcounter{utterance}  
    & \multicolumn{4}{p{0.6\linewidth}}{
        \cellcolor[rgb]{0.9,0.9,0.9}{
            \makecell[{{p{\linewidth}}}]{
                \texttt{\tiny{[P1$\rangle$GM]}}
                \texttt{СОГЛАСЕН: \{'C08', 'C43', 'C73', 'B19', 'C10', 'C61', 'C55'\}} \\
            }
        }
    }
    & & \\ \\

    \theutterance \stepcounter{utterance}  
    & & \multicolumn{4}{p{0.6\linewidth}}{
        \cellcolor[rgb]{0.9,0.9,0.9}{
            \makecell[{{p{\linewidth}}}]{
                \texttt{\tiny{[P1$\langle$GM]}}
                \texttt{Ваш ответ не содержал аргумента. Вы должны включить структурированный формат 'АРГУМЕНТ: \{...\}' где{-}то в вашем ответе. Попробуйте снова} \\
            }
        }
    }
    & \\ \\

    \theutterance \stepcounter{utterance}  
    & \multicolumn{4}{p{0.6\linewidth}}{
        \cellcolor[rgb]{0.9,0.9,0.9}{
            \makecell[{{p{\linewidth}}}]{
                \texttt{\tiny{[P1$\rangle$GM]}}
                \texttt{СОГЛАСЕН: \{'C08', 'C43', 'C73', 'B19', 'C10', 'C61', 'C55'\}} \\
\texttt{АРГУМЕНТ: \{'Я подтверждаю своё согласие с данным набором, так как он валиден по бюджету и максимизирует мою выгоду.'\}} \\
            }
        }
    }
    & & \\ \\

    \theutterance \stepcounter{utterance}  
    & & & \multicolumn{2}{p{0.3\linewidth}}{
        \cellcolor[rgb]{0.9,0.9,0.9}{
            \makecell[{{p{\linewidth}}}]{
                \texttt{\tiny{[GM$|$GM]}}
                \texttt{abort} \\
            }
        }
    }
    & & \\ \\

    \theutterance \stepcounter{utterance}  
    & & & \multicolumn{2}{p{0.3\linewidth}}{
        \cellcolor[rgb]{0.9,0.9,0.9}{
            \makecell[{{p{\linewidth}}}]{
                \texttt{\tiny{[GM$|$GM]}}
                \texttt{end game} \\
            }
        }
    }
    & & \\ \\

\end{supertabular}
}

\end{document}
